\begin{definition}[Infinite series and convergence of a series] \label{definition:infinite_series_and_convergence_of_infinite_series_in_a_topological_abelian_semigroup}
    Let $A$ be a \CrefAndHyperrefIfExist{definition:topological_semigroup}{topological abelian semigroup} written additively. Let $(a_n)_{n=1}^\infty$ be a sequence in $A$. The formal expression
    \hl{$\sum_{n=1}^\infty a_n$}
    is called an \hldef{infinite series}.

    Let $(S_N)_{N=1}^\infty$ be the sequence of partial sums of $(a_n)$, where $S_N = \sum_{n=1}^N a_n$. The infinite series $\sum_{n=1}^\infty a_n$ is said to \hldef{converge} if there exists an element $S \in A$ such that for every open neighborhood $U$ of $S$ in the topological space $A$\CrefIfExists{definition:topological_space}, there exists $N_0 \geq 1$ such that $S_N \in U$ for all $N \geq N_0$, i.e., if the sequence $(S_N)$ \CrefAndHyperrefIfExist{definition:convergence_and_limit_of_sequence_in_a_topological_space}{converges to} $S$ in the topology of $A$. If such an $S$ exists, then $S$ is called the \hldef{sum of the infinite series $\sum_{n=1}^\infty a_n$}, and we write
    $$\hlin{\sum_{n=1}^\infty a_n = S.}$$
    If $(S_N)$ does not converge, then the series is said to \hldef{diverge}.

    If $A$ is a topological abelian \CrefAndHyperrefIfExist{definition:monoid}{monoid} (i.e., it has an additive identity $0_A$), then by convention the empty sum over no indices is taken to be $0_A$. This convention is used, for example, to define the $0$th partial sum $S_0 = 0_A$ when convenient.

    If the semigroup operation is written multiplicatively rather than additively, the same definition with $\sum$ replaced by $\prod$ and $S_N$ replaced by the partial products $P_N = \prod_{n=1}^N a_n$ defines an \hldef{infinite product}, denoted
    $$\hlin{\prod_{n=1}^\infty a_n,}$$
    whose limit (when it exists) is called the \hldef{value of the infinite product}. If $A$ is a topological abelian monoid with identity $1_A$, the empty product is taken to be $1_A$ by convention.

    \TextIfExists{proposition:convergence_of_a_monotone_sequence_to_a_supremum_or_infimum_in_a_totally_ordered_set_coincides_with_topological_convergence_in_the_order_topology}{
        When $A$ is a \CrefAndHyperrefIfExist{definition:ordered_abelian_group}{totally ordered abelian group} equipped with the \CrefAndHyperrefIfExist{definition:order_topology_on_a_partially_ordered_set}{order topology} and the terms $a_n$ are non-negative, the sequence of partial sums $(S_N)$ is nondecreasing, so this topological notion of convergence coincides with the order-theoretic notion of \CrefAndHyperrefIfExist{definition:convergence_of_a_monotone_sequence_in_a_partially_ordered_set_with_adjoined_infinity}{convergence of a monotone sequence to a supremum} whenever the relevant supremum exists.
    }
\end{definition}